Les critères que notre plateforme doit satisfaire sont : 
\begin{itemize}[label=$\blacktriangleright$]
\item \textbf{Ergonomie :}

L'ergonomie sera une priorité essentielle pour garantir une expérience utilisateur optimale. Nous développerons l'interface avec \textbf{shadcn/ui} et \textbf{Tailwind CSS}, assurant une cohérence visuelle et une utilisabilité maximale.
\newline
Nous intégrerons un système de thématisation dynamique permettant de basculer entre mode clair et mode sombre. Cette fonctionnalité réduira la fatigue oculaire lors d'utilisations prolongées et s'adaptera aux préférences individuelles des utilisateurs.
\newline
Nous implémenterons une architecture \textit{responsive} via les \textit{breakpoints} de \textbf{Tailwind CSS}, garantissant une adaptation automatique de l'interface aux différentes résolutions d'écran (ordinateur de bureau, tablette, mobile). Les composants seront conçus pour assurer une expérience fluide sur tous les supports.
\newline
Nous exploiterons les composants réactifs basés sur \textbf{React} pour permettre des mises à jour en temps réel sans rechargement de page. Les transitions et animations seront subtiles et fonctionnelles, guidant l'utilisateur de manière intuitive sans le distraire de ses tâches principales.
\newline
Par ailleurs, un système de retour d'information (\textit{feedback}) sera mis en place pour informer l'utilisateur du résultat de ses actions via des messages de confirmation (« Ticket créé avec succès ! », « Votre profil a été mis à jour. », « Votre message a bien été envoyé ! »), garantissant ainsi une communication claire de l'état du système.

\item \textbf{Sécurité :}

La sécurité sera une priorité absolue pour notre application. Nous intégrerons des mesures robustes telles que le hachage des mots de passe avec \textbf{bcrypt} (facteur de coût 10) et l'utilisation de jetons d'accès JWT (\textit{JSON Web Token}) pour l'authentification.
\newline
Nous implémenterons un système de limitation de débit (\textit{rate limiting}) pour prévenir les attaques par force brute (limite de 5 tentatives de connexion par 15 minutes) et les abus d'API (100 requêtes par minute par adresse IP). Un mécanisme de contrôle d'accès reposant sur le principe \textbf{RBAC} restreindra les accès selon les droits attribués, renforçant la protection des données sensibles.
\newline
Afin de garantir une authentification sécurisée, l'application prendra en charge la connexion via Google et GitHub grâce au protocole OAuth 2.0.
\newline
Par ailleurs, la génération d'un jeton de réinitialisation de mot de passe sera associée à une durée de validité de 24 heures et à un usage unique. Afin d'assurer un audit complet, notre application garantira la journalisation de toutes les actions avec les informations nécessaires pour identifier l'acteur ainsi que l'action réalisée.
\newline
Nous implémenterons des protections contre les attaques courantes : jetons \textbf{CSRF} pour prévenir les requêtes intersites malveillantes, validation et échappement des entrées utilisateur contre les injections \textbf{XSS}, et requêtes paramétrées via \textbf{Prisma ORM} contre les injections SQL.
\newline
Lors de la création d'un contrat avec l'option de téléversement, chaque fichier téléversé sera vérifié à l'aide d'un \textit{hash} SHA-256, permettant de détecter les doublons et d'éviter les téléversements multiples d'un même contrat, garantissant ainsi l'unicité contractuelle.
\newline
Enfin, lors de la création d'un utilisateur, le système générera automatiquement un mot de passe temporaire sécurisé. Par conséquent, lors de sa première connexion, l'utilisateur sera contraint de le modifier afin de renforcer la sécurité de son compte.

\item \textbf{Performance :}

Les performances de notre plateforme seront optimisées à plusieurs niveaux pour garantir des temps de réponse rapides et une expérience utilisateur fluide.
\newline
% À VÉRIFIER : L'architecture sera conçue pour supporter une charge de 1000 utilisateurs simultanés avec possibilité d'évolution horizontale via l'ajout de serveurs supplémentaires. Cette approche garantira la scalabilité de la plateforme en fonction de la croissance du nombre d'utilisateurs et du volume de données.
% \newline
L'architecture sera conçue avec une approche modulaire permettant une évolution horizontale via l'ajout de serveurs supplémentaires si nécessaire. Cette approche garantira la scalabilité de la plateforme en fonction de la croissance du nombre d'utilisateurs et du volume de données.
\newline
Au niveau de la base de données, nous implémenterons une stratégie d'indexation appropriée sur les colonnes fréquemment interrogées (identifiants, dates, statuts, clés étrangères). L'utilisation de \textbf{Prisma ORM} permettra de générer des requêtes SQL optimisées et d'exploiter les fonctionnalités avancées de PostgreSQL telles que les index composites et les index partiels pour les requêtes complexes.
\newline
Afin d'éviter les problèmes de requêtes N+1, les relations entre entités seront chargées de manière sélective via les mécanismes \textbf{select} et \textbf{include} de Prisma. Les requêtes complexes utiliseront le chargement anticipé (\textit{eager loading}) pour minimiser le nombre d'appels à la base de données.
\newline
Enfin, nous implémenterons la division du code (\textit{code splitting}) et le chargement différé (\textit{lazy loading}) des composants pour réduire la taille du paquet JavaScript initial. Les requêtes API seront optimisées via le préchargement (\textit{prefetching}) et la mise en cache côté client, minimisant les appels réseau redondants.

% \item \textbf{Disponibilité :}
% 
% La disponibilité sera une exigence critique pour notre plateforme de gestion de maintenance. Nous garantirons une disponibilité de \textbf{99.5\%} (temps de fonctionnement), correspondant à un temps d'arrêt maximal de 3.65 heures par mois, ainsi qu'une surveillance continue des délais de réponse et de résolution des tickets conformément aux engagements contractuels SLA.
% \newline
% Nous implémenterons la communication en temps réel via le protocole \textbf{WebSocket}, permettant des échanges bidirectionnels instantanés entre le serveur et les clients. Cette technologie supportera la messagerie instantanée entre l'équipe et les clients, ainsi que la diffusion des notifications en temps réel. En cas d'indisponibilité WebSocket (pare-feu restrictif), un mécanisme de repli via interrogation HTTP (\textit{polling}) sera automatiquement activé pour maintenir la communication.
% \newline
% Afin de garantir qu'aucune information ne soit perdue, nous mettrons en place un mécanisme de persistance robuste. Les notifications seront enregistrées en base de données avec un système de synchronisation automatique. En cas de déconnexion temporaire, les messages en attente seront automatiquement livrés lors de la reconnexion de l'utilisateur.
% \newline
% Par ailleurs, des sauvegardes automatiques de la base de données seront effectuées quotidiennement avec une rétention de 30 jours. Un plan de reprise d'activité (PRA) permettra une restauration en moins de 4 heures (RTO - \textit{Recovery Time Objective}) avec une perte de données maximale de 1 heure (RPO - \textit{Recovery Point Objective}).
% \newline
% Un système de surveillance (\textit{monitoring}) en temps réel détectera automatiquement les anomalies de performance et de disponibilité, déclenchant des alertes vers l'équipe technique. Les métriques surveillées incluront l'utilisation CPU, mémoire, espace disque, temps de réponse des requêtes et taux d'erreurs.
% \newline
% Un système d'alertes proactives sera mis en place pour prévenir l'expiration des contrats. Les notifications seront déclenchées à trois intervalles critiques : 30 jours, 15 jours et 7 jours avant l'échéance. Ces alertes seront envoyées par courriel et affichées dans l'interface, garantissant que les responsables disposent d'un délai suffisant pour renouveler les contrats.
% \newline
% Enfin, le système de surveillance des SLA calculera automatiquement les temps de réponse et de résolution, déclenchant des alertes en cas de risque de dépassement. Les tableaux de bord afficheront en temps réel les indicateurs de performance, permettant une intervention proactive avant la violation des engagements contractuels.

\item \textbf{Maintenabilité :}

La maintenabilité sera assurée par l'adoption de bonnes pratiques de développement et une architecture modulaire facilitant l'évolution du système. Le code source sera structuré selon les principes SOLID et les conventions de codage de chaque technologie (\textbf{Next.js}, \textbf{NestJS}), garantissant la lisibilité et la facilité de modification.
\newline
Nous implémenterons une documentation technique complète pour les API \textit{backend} générée automatiquement via \textbf{Swagger/OpenAPI}. Cette documentation sera maintenue à jour automatiquement lors de chaque modification du code, assurant sa cohérence avec l'implémentation.
\newline
Enfin, afin de garantir la qualité du code, nous viserons une couverture de tests minimale de 80\% : tests unitaires avec \textbf{Jest} pour les fonctions métier, tests d'intégration pour les API \textit{backend}, et tests end-to-end avec \textbf{Playwright} pour les parcours utilisateur critiques. Ces tests seront exécutés automatiquement avant chaque déploiement.
\end{itemize}
